\documentclass{exam}

\usepackage[margin=1in]{geometry}

\usepackage{comment}
\usepackage{graphicx}
\usepackage{hyperref}
\usepackage{listings}

\pagestyle{headandfoot}
\firstpageheader{CSCI 221}{}{Assembly Basics Assignment}
\runningheader{CSCI 221}{}{Assembly Basics Assignment}


\begin{document}
\title{Assembly Basics Assignment (Evaluation)}
\author{CSCI 221: Computer Science Fundamentals 2}
\date{Spring 2026}
\maketitle


This assignment is an opportunity to evaluate your understanding of assembly basics. 
Feel free to collaborate with others and use resources, but all code and writeups must be your own. 
To submit your assignment, please submit your code and writeup. 
Mark which sections of the code are for which questions by using different files or comments in the file. 
For this assignment, your code does not need to check for invalid input. 
For this assignment, you should not use pseudoinstructions. 

\textbf{Due Date:}
Friday, April 10th at 1:00 pm. 

For this assignment, you will be writing code in MIPS. 
Unfortunately, most machines you interact with will use a different ISA than MIPS, which means that they will be unable to run MIPS code. 
In order to test and debug your code, you will need to make use of a MIPS simulator. 
Many MIPS simulators are available online. 
If you have a preferred simulator, feel free to use it. 
If not, we recommend the QtSPIM MIPS simulator. 
If you are having issues installing on a mac, you can try the following link: \url{https://support.apple.com/guide/mac-help/open-a-mac-app-from-an-unknown-developer-mh40616/mac}. 

\begin{questions}

\question[10]
\textbf{Fibonacci.}
Write MIPS assembly code that computes the $n$th Fibonacci number. 
For reference, the 0th Fibonacci number is 0, the 1st Fibonacci number is 1, and every subsequent Fibonacci number is the sum of the two previous Fibonacci numbers. 
You should use loops rather than recursion. 
You should store the result in $\$v0$. 
Assume that $n$ is stored in $\$a0$. 

\question[10]
\textbf{Array Counting.}
Write MIPS assembly code that counts the frequency of values in particular ranges in a given array of single-precision floating point numbers. 
In addition to the array, you are given 2 values. 
You should count the number of values in the array smaller than both values, larger than both values, and in between the values. 
Assume that the lower end of each range is inclusive, and the upper end of each range is exclusive. 
You should store the results of the counting in $\$s0$, $\$s1$, and $\$s2$. 
Assume that the array is stored in $\$a0$ and its length is stored in $\$a1$. 
Assume that the 2 values are stored in  $\$a2$, and $\$a3$, and that they are stored in ascending order. 

\question[10]
\textbf{Array Permutation.}
Write MIPS assembly code that performs the same functionality as the following C code:
\begin{verbatim}
for (int i=0; i<length; i++){
    swap(&array[i], &array[permutations[i]]);
}
\end{verbatim}
Assume that \texttt{array} and \texttt{permutations} both consist of 32-bit integers and that each value in \texttt{permutations} is between 0 and \texttt{length}. 
Assume that the arrays contain the same number of elements. 
Assume that \texttt{array}, \texttt{permutations}, and \texttt{length} are stored in $\$a0$, $\$a1$, and $\$a2$ respectively. 

\question[10]
\textbf{Linked List Product.}
Write MIPS assembly code that computes the product of all values in a linked list. 
Assume that nodes in the list consist of a 16-bit integer and a pointer to the next node, stored in that order. 
Assume that the final node's pointer contains the value 0. 
Assume that the result will fit in 64 bits. 
Remember that pointers in MIPS are 32 bits. 
You should store the result in $\$v0$ and $\$v1$. 
Assume that a pointer to the first node in the list is stored in $\$a0$. 

%\question[10]
%\textbf{Multiplying Arrays.}
%\begin{parts}
%\part[4]
%Write a MIPS assembly function that takes as input two 32 bit integers and returns their product. 
%
%\part[6]
%Write a MIPS assembly function that takes as input pointers to two lists of 32 bit integers, a pointer to a list of 64 bit integers and the length of those lists, which is assumed to be the same. 
%(All lists contain the same number of elements.)
%Each element in the third list should be set to the product of the corresponding elements in the first two lists. 
%In code form, this means:
%\begin{verbatim}
%for (int i=0; i<length; i++){
%    array3[i] = array1[i] * array2[i];
%}
%\end{verbatim}
%For full credit, use your function from the previous part. 
%\end{parts}


\question[10]
\textbf{Writeup.}
In addition to your code, you should submit a brief (~1-2 pages) writeup that describes the state of your code. 
You should discuss what is working, what is not, and any known errors. 
In this discussion, you should reference the tests you did and how they support your claims. 
You will need to provide a brief description of your tests for you to reference. 
Your writeup should be in a separate .txt or .pdf file. 

\end{questions}

\end{document}